
\documentclass{article}
\usepackage{colortbl}
\usepackage{makecell}
\usepackage{multirow}
\usepackage{supertabular}

\begin{document}

\newcounter{utterance}

\centering \large Interaction Transcript for game `wordle', experiment `medium\_frequency\_words\_no\_clue\_no\_critic', episode 5 with gemma{-}4{-}26b{-}a4b{-}it.
\vspace{24pt}

{ \footnotesize  \setcounter{utterance}{1}
\setlength{\tabcolsep}{0pt}
\begin{supertabular}{c@{$\;$}|p{.15\linewidth}@{}p{.15\linewidth}p{.15\linewidth}p{.15\linewidth}p{.15\linewidth}p{.15\linewidth}}
    \# & \multicolumn{2}{c}{Player} && \multicolumn{2}{c}{Game Master} \\
    \hline

    \theutterance \stepcounter{utterance}  
    & & & \multicolumn{4}{p{0.6\linewidth}}{
        \cellcolor[rgb]{0.9,0.9,0.9}{
            \makecell[{{p{\linewidth}}}]{
                \texttt{\tiny{[P1$\langle$GM]}}
                \texttt{Vous êtes un magicien du langage qui aime deviner des mots en utilisant les règles données.} \\
\\ 
\texttt{Bienvenue à Wordle ! Vous avez six tentatives pour deviner le mot cible, un mot valide de cinq lettres minuscules. Veuillez utiliser les balises "explication :" et "proposition :" pour fournir une explication concise pour chaque proposition.} \\
\\ 
\texttt{Par exemple, si votre proposition est "année", votre réponse devrait être} \\
\texttt{explication : c'est un mot commun de cinq lettres, et je commence ma proposition avec ce mot.} \\
\texttt{proposition : année} \\
\\ 
\texttt{Après chaque proposition, votre réponse sera validée, et vous recevrez un retour indiquant quelles lettres sont correctes (vert), quelles lettres sont correctes mais à la mauvaise position (jaune), et quelles lettres sont incorrectes (rouge). Ce retour peut être utile pour déterminer quelles lettres inclure ou exclure dans votre prochaine proposition.} \\
\\ 
\texttt{Par exemple, le retour pour "année" pourrait être :} \\
\texttt{retour\_proposition : a<jaune> n<jaune> n<vert> é<jaune> e<rouge>} \\
\\ 
\texttt{L'explication doit contenir des détails sur la manière dont le retour\_proposition est utilisé pour arriver à une nouvelle proposition.} \\
\\ 
\texttt{Ne générez aucun autre texte.} \\
\\ 
\texttt{Commençons avec votre première proposition.} \\
            }
        }
    }
     \\ \\

    \theutterance \stepcounter{utterance}  
    & \multicolumn{4}{p{0.6\linewidth}}{
        \cellcolor[rgb]{0.561,0.737,0.561}{
            \makecell[{{p{\linewidth}}}]{
                \texttt{\tiny{[P1$\rangle$GM]}}
                \texttt{explication : je commence par un mot courant contenant des voyelles fréquentes pour tester les bases du mot cible.} \\
\texttt{proposition : adieu} \\
            }
        }
    }
    & & \\ \\

    \theutterance \stepcounter{utterance}  
    & & & \multicolumn{2}{p{0.3\linewidth}}{
        \cellcolor[rgb]{0.9,0.9,0.9}{
            \makecell[{{p{\linewidth}}}]{
                \texttt{\tiny{[GM$|$GM]}}
                \texttt{attempts = 1} \\
\texttt{target\_word = label} \\
\texttt{guess = adieu} \\
\texttt{guess\_feedback = a<yellow> d<red> i<red> e<green> u<red>} \\
            }
        }
    }
    & & \\ \\

    \theutterance \stepcounter{utterance}  
    & & & \multicolumn{4}{p{0.6\linewidth}}{
        \cellcolor[rgb]{0.9,0.9,0.9}{
            \makecell[{{p{\linewidth}}}]{
                \texttt{\tiny{[P1$\langle$GM]}}
                \texttt{retour\_proposition : a<jaune> d<rouge> i<rouge> e<vert> u<rouge>} \\
\\ 
\texttt{Fournissez votre réponse uniquement dans ce format.} \\
\texttt{explication : détails} \\
\texttt{proposition : mot} \\
            }
        }
    }
     \\ \\

    \theutterance \stepcounter{utterance}  
    & \multicolumn{4}{p{0.6\linewidth}}{
        \cellcolor[rgb]{0.561,0.737,0.561}{
            \makecell[{{p{\linewidth}}}]{
                \texttt{\tiny{[P1$\rangle$GM]}}
                \texttt{explication : la lettre 'e' est à la bonne place (position 4), la lettre 'a' est présente mais pas en position 1, et les lettres 'd', 'i', 'u' sont exclues. Je cherche un nouveau mot avec 'e' en 4ème position et incluant 'a'.} \\
\texttt{proposition : table} \\
            }
        }
    }
    & & \\ \\

    \theutterance \stepcounter{utterance}  
    & & & \multicolumn{2}{p{0.3\linewidth}}{
        \cellcolor[rgb]{0.9,0.9,0.9}{
            \makecell[{{p{\linewidth}}}]{
                \texttt{\tiny{[GM$|$GM]}}
                \texttt{attempts = 2} \\
\texttt{target\_word = label} \\
\texttt{guess = table} \\
\texttt{guess\_feedback = t<red> a<green> b<green> l<yellow> e<yellow>} \\
            }
        }
    }
    & & \\ \\

    \theutterance \stepcounter{utterance}  
    & & & \multicolumn{4}{p{0.6\linewidth}}{
        \cellcolor[rgb]{0.9,0.9,0.9}{
            \makecell[{{p{\linewidth}}}]{
                \texttt{\tiny{[P1$\langle$GM]}}
                \texttt{retour\_proposition : t<rouge> a<vert> b<vert> l<jaune> e<jaune>} \\
\\ 
\texttt{Fournissez votre réponse uniquement dans ce format.} \\
\texttt{explication : détails} \\
\texttt{proposition : mot} \\
            }
        }
    }
     \\ \\

    \theutterance \stepcounter{utterance}  
    & \multicolumn{4}{p{0.6\linewidth}}{
        \cellcolor[rgb]{0.561,0.737,0.561}{
            \makecell[{{p{\linewidth}}}]{
                \texttt{\tiny{[P1$\rangle$GM]}}
                \texttt{explication : la lettre 'a' est en position 2 et 'b' est en position 3. La lettre 'l' est présente mais pas en position 4, et 'e' est présente mais pas en position 5. Les lettres 't', 'd', 'i', 'u' sont exclues. Je teste une nouvelle combinaison avec 'a' et 'b' fixes.} \\
\texttt{proposition : sable} \\
            }
        }
    }
    & & \\ \\

    \theutterance \stepcounter{utterance}  
    & & & \multicolumn{2}{p{0.3\linewidth}}{
        \cellcolor[rgb]{0.9,0.9,0.9}{
            \makecell[{{p{\linewidth}}}]{
                \texttt{\tiny{[GM$|$GM]}}
                \texttt{attempts = 3} \\
\texttt{target\_word = label} \\
\texttt{guess = sable} \\
\texttt{guess\_feedback = s<red> a<green> b<green> l<yellow> e<yellow>} \\
            }
        }
    }
    & & \\ \\

    \theutterance \stepcounter{utterance}  
    & & & \multicolumn{4}{p{0.6\linewidth}}{
        \cellcolor[rgb]{0.9,0.9,0.9}{
            \makecell[{{p{\linewidth}}}]{
                \texttt{\tiny{[P1$\langle$GM]}}
                \texttt{retour\_proposition : s<rouge> a<vert> b<vert> l<jaune> e<jaune>} \\
\\ 
\texttt{Fournissez votre réponse uniquement dans ce format.} \\
\texttt{explication : détails} \\
\texttt{proposition : mot} \\
            }
        }
    }
     \\ \\

    \theutterance \stepcounter{utterance}  
    & \multicolumn{4}{p{0.6\linewidth}}{
        \cellcolor[rgb]{0.561,0.737,0.561}{
            \makecell[{{p{\linewidth}}}]{
                \texttt{\tiny{[P1$\rangle$GM]}}
                \texttt{explication : 'a' est en position 2, 'b' est en position 3. 'l' et 'e' sont présentes mais pas aux positions testées (4 et 5). 's' est exclu. Je tente de placer 'l' et 'e' ailleurs, par exemple 'l' en position 1 et 'e' en position 4.} \\
\texttt{proposition : label} \\
            }
        }
    }
    & & \\ \\

    \theutterance \stepcounter{utterance}  
    & & & \multicolumn{2}{p{0.3\linewidth}}{
        \cellcolor[rgb]{0.9,0.9,0.9}{
            \makecell[{{p{\linewidth}}}]{
                \texttt{\tiny{[GM$|$GM]}}
                \texttt{attempts = 4} \\
\texttt{target\_word = label} \\
\texttt{guess = label} \\
\texttt{guess\_feedback = l<green> a<green> b<green> e<green> l<green>} \\
            }
        }
    }
    & & \\ \\

    \theutterance \stepcounter{utterance}  
    & & & \multicolumn{2}{p{0.3\linewidth}}{
        \cellcolor[rgb]{0.9,0.9,0.9}{
            \makecell[{{p{\linewidth}}}]{
                \texttt{\tiny{[GM$|$GM]}}
                \texttt{game\_result = WIN} \\
            }
        }
    }
    & & \\ \\

\end{supertabular}
}

\end{document}
